% =====================================================================================
% Step 3 output (docs/article/PAPER_EXECUTION_PLAN.md). Opus 5 pass, 2026-09-17.
%
% Front half (pseudocode extraction from code, Sonnet 5) is paper/equations_pseudocode.md.
% This file is the LaTeX definition of record. Rule A5: every equation carries an
% `equation_operation` line stating what the equation ACTUALLY computes, derived from the
% executed code rather than from a docstring or a textbook formulation, plus the file and
% line it was read off. Line numbers verified against the working tree on 2026-09-17.
%
% Where the code and its own documentation disagree, the code wins and the disagreement is
% recorded in the note under the equation. Four such gaps survive into the write-up:
%   E2  the docstring's left-multiplication order is not the executed tensor layout;
%   E3  the intercept column's ridge penalty is zeroed, which no comment near the solver says;
%   E4  the "joint loss" is descended on two different optimizer schedules, so no single
%       scalar is literally minimized;
%   E7  C(I+A)(B_1+B_2 y) is a special case, not the controller's decision rule.
%
% This file defines no macros and loads no packages; main.tex already provides amsmath and
% amssymb. It is not \input by main.tex -- it is the source the Method section draws from.
% =====================================================================================

% -------------------------------------------------------------------------------------
% E1. Delay-embedded state construction (Markov / Memory-L / Augmented-L)
% code: src/koopman_ae/core.py:169 (_state_at_turn), :180 (build_augmented_state_dataset)
%       scripts/train.py:98 (_state_config) selects the variant
% equation_operation: concatenate a fixed number of lagged readout values, and for the
%   augmented family an equal number of lagged control values formed as the tracking error
%   r - y_t, into one delay-embedded state vector; the number of terms is set per variant
%   outside the state builder (markov = 1 term regardless of lag, memory-L and augmented-L
%   = L+1 terms).
% -------------------------------------------------------------------------------------
\begin{align}
  u_t &= \phi\!\left(r - y_t\right), &
  \phi(e) &=
  \begin{cases}
    e, & \text{control\_mode} = \texttt{error},\\[2pt]
    \bigl[\,e,\ \lvert e\rvert,\ \operatorname{sign}(e)\,\bigr], & \text{control\_mode} = \texttt{error\_abs\_sign},
  \end{cases}
  \label{eq:control-from-error}\\[4pt]
  z_t &= \bigl[\,y_t,\ y_{t-1},\ \dots,\ y_{t-(L_y-1)},\ u_t,\ u_{t-1},\ \dots,\ u_{t-(L_u-1)}\,\bigr],
  \label{eq:delay-state}
\end{align}
%
\noindent with $(L_y, L_u) = (1,0)$ for \textsc{markov}, $(L+1, 0)$ for \textsc{memory-}$L$,
and $(L+1, L+1)$ for \textsc{augmented-}$L$. A row enters the design matrix only when
$t \ge \max(L_y, L_u)$, when turn $t{+}1$ exists, and when every lagged turn the state needs
is present.

\paragraph{Note (structural, carried into Limitations).}
On the \texttt{core} tasks $r$ is one fixed vector per trajectory, so by
\eqref{eq:control-from-error} the control $u_t$ is an exact affine function of the state's
own leading block $y_t$. The design matrix $[Z \mid R \mid 1]$ is therefore rank deficient on
all eight \texttt{core} tasks ($R^2 = 1.0$ recovering $u_t$ from $z_t$), and no quantity that
requires separating the action's effect from the state's is identifiable on that line.

% -------------------------------------------------------------------------------------
% E2. Autoencoder lift and latent linear dynamics
% code: src/koopman_ae/core.py:514 (_latent_step_tensor), :999 (encode), :1005 (decode),
%       :1011 (predict_next_latent), :1021 (predict_next_z)
% equation_operation: encode the delay-embedded state into a latent vector, advance that
%   latent vector one step by an affine map driven by the per-trajectory target vector r,
%   decode back to state space, and read the observable off a FIXED slice of the decoded
%   state -- the readout is a projection, not a learned head.
% -------------------------------------------------------------------------------------
\begin{align}
  \xi_t &= \mathcal{E}_\theta(z_t), &
  \xi_{t+1} &= K\,\xi_t + B\,r + c, &
  \hat z_{t+1} &= \mathcal{D}_\theta(\xi_{t+1}), &
  \hat y_{t+1} &= \bigl[\hat z_{t+1}\bigr]_{1:d_y}.
  \label{eq:ae-koopman}
\end{align}
%
\noindent A horizon-$H$ rollout iterates the middle map $H$ times from a single encode, with
$r$ held fixed throughout, and decodes at every step.

\paragraph{Note (code vs.\ docstring).}
The class docstring writes \eqref{eq:ae-koopman} in the left-multiplication order above; the
executed line is \texttt{xi @ K.T + r @ B.T + c}, a batched row-vector convention. The two are
the same map once transposes are accounted for. \eqref{eq:ae-koopman} is written in the
column-vector convention because that is what the surrounding prose uses; any shape claim in
the text must be stated for that convention and not read off the tensor code.

% -------------------------------------------------------------------------------------
% E3. Two-stage estimator: reconstruction by gradient descent, then closed-form ridge
% code: src/koopman_ae/core.py:39 (_ridge_solve), :712 (_fit_latent_ridge),
%       :730 (fit, the training_mode == "reconstruction_then_ridge" branch)
% equation_operation: fit the encoder and decoder by gradient descent on reconstruction loss
%   ALONE, then freeze the encoder and solve one ridge least-squares problem in closed form
%   for (K, B, c) against the augmented design [xi_t, r, 1], with the appended constant
%   column's own ridge coefficient forced to zero so the intercept is unpenalized while K and
%   B are shrunk.
% -------------------------------------------------------------------------------------
\begin{align}
  \theta^{\mathrm{ae}} &= \arg\min_{\theta}\ \lambda_{\mathrm{rec}}\,
    \bigl\lVert \mathcal{D}_\theta(\mathcal{E}_\theta(z)) - z \bigr\rVert_2^2,
  \label{eq:stage-one}\\[4pt]
  X &= \bigl[\ \Xi_t \ \big|\ R \ \big|\ \mathbf{1}\ \bigr], \qquad
  \Lambda = \alpha\,\mathrm{diag}(1,\dots,1,0),
  \label{eq:stage-two-design}\\[4pt]
  \bigl[\,K^\top;\ B^\top;\ c^\top\,\bigr]
    &= \bigl(X^\top X + \Lambda\bigr)^{+} X^\top\, \Xi_{t+1}.
  \label{eq:stage-two-solve}
\end{align}
%
\noindent $\Xi_t = \mathcal{E}_{\theta^{\mathrm{ae}}}(Z_t)$ is a no-grad forward pass through
the frozen encoder, $(\cdot)^{+}$ is the pseudoinverse, and the final $0$ on $\Lambda$'s
diagonal is the unregularized intercept.

\paragraph{Note (undocumented in code).}
The zeroed intercept penalty, $\Lambda_{dd}=0$, appears nowhere in the solver's comments or
docstrings. It matters: with a penalized intercept the fitted dynamics absorb the latent mean
into $K$, and the estimator is no longer the one described here. State it in the Method text
rather than leaving a reader to infer it.

% -------------------------------------------------------------------------------------
% E4. The joint objective's four terms, and the two schedules that descend them
% code: src/koopman_ae/core.py:845-892 (per-batch loss and the separate per-epoch step),
%       :977 (_multi_step_loss), :517 (_eval_loss restates the same arithmetic no-grad)
% equation_operation: weight four mean-squared-error terms -- state reconstruction, one-step
%   prediction in state space, one-step prediction in latent space, and a multi-step latent
%   rollout -- but descend them on TWO schedules: the first three take one gradient step per
%   minibatch, the fourth takes one additional gradient step per epoch. No single scalar is
%   minimized by a single step; the logged "loss" is the weighted sum for early stopping only.
% -------------------------------------------------------------------------------------
\begin{align}
  \mathcal{L}_{\mathrm{batch}}
    &= \lambda_{\mathrm{rec}} \bigl\lVert \mathcal{D}(\mathcal{E}(z_t)) - z_t \bigr\rVert^2
     + \lambda_{\mathrm{pred}} \bigl\lVert \mathcal{D}(K\mathcal{E}(z_t)+Br+c) - z_{t+1} \bigr\rVert^2
     \nonumber\\
    &\qquad + \lambda_{\mathrm{lat}} \bigl\lVert K\mathcal{E}(z_t)+Br+c - \mathcal{E}(z_{t+1}) \bigr\rVert^2,
  \label{eq:joint-batch}\\[4pt]
  \mathcal{L}_{\mathrm{multi}}
    &= \lambda_{\mathrm{multi}}\,
       \frac{1}{|\mathcal{S}|}\sum_{s \in \mathcal{S}} \frac{1}{H_s}\sum_{k=1}^{H_s}
       \bigl\lVert \mathcal{D}\bigl(\Phi^k(\mathcal{E}(z_0^{s}), r^{s})\bigr) - z_k^{s} \bigr\rVert^2,
  \label{eq:joint-multi}
\end{align}
%
\noindent where $\Phi^k$ is $k$ applications of the affine latent map of
\eqref{eq:ae-koopman}. Each epoch takes $n_{\mathrm{batches}}$ gradient steps on
\eqref{eq:joint-batch} and exactly one on \eqref{eq:joint-multi}.

\paragraph{Note (why this is not one objective).}
$\lambda_{\mathrm{multi}}$ is therefore not the only knob controlling the multi-step term's
influence: that term receives $1/(n_{\mathrm{batches}}+1)$ of the epoch's steps, and
$n_{\mathrm{batches}}$ scales with dataset size. The same $\lambda_{\mathrm{multi}}$ means
different things on datasets of different size, so a cross-task sweep over it confounds
$\lambda$ with $n$. This is the stated reason the multi-step ablation is not reported: the
paper claims only that the multi-step term was never swept at comparable strength.

% -------------------------------------------------------------------------------------
% E5. Behavioral-line state schema and its lifted linear dynamics
% code: persona_drift_control/src/persona_drift/modeling/dataset.py:33 (ReducedStateConfig),
%       :163 (build_reduced_state_pairs);
%       .../modeling/koopman.py:105 (_psi), :114 (fit), :139 (step), :146 (readout)
% equation_operation: build a delay-embedded reduced state from nu readout lags (including
%   the current one) and mu control lags (excluding the current slot), lift it through an
%   optional feature dictionary, and fit an affine one-step map plus a SEPARATE ridge readout
%   by ridge normal equations -- two independent least-squares problems, not one joint fit.
% -------------------------------------------------------------------------------------
\begin{align}
  z_t &= \bigl[\,y_{t-\nu+1},\dots,y_t,\ v_{t'-\mu},\dots,v_{t'-1},\ a_t\,\bigr],
        \qquad t' = t + \mathbb{1}[\text{contemporaneous}],
  \label{eq:beh-state}\\[4pt]
  \eta_t &= \psi(z_t) = \bigl[\,z_t,\ \varphi(z_t)\,\bigr],
  \label{eq:beh-lift}\\[4pt]
  \eta_{t+1} &= A\,\eta_t + B\,v_t + b,
  \qquad
  \hat y_t = C\,\eta_t.
  \label{eq:beh-dynamics}
\end{align}
%
\noindent $a_t$ collects auxiliary columns when configured, and $\varphi$ is the empty map for
the ARX variant. $(A,B,b)$ solve
$\min \lVert [\Psi \mid V \mid \mathbf{1}]\,\Theta - \Psi_{+}\rVert^2 + \rho\lVert\Theta\rVert^2$;
$C$ solves its own ridge problem $\min \lVert \Psi C^\top - Y\rVert^2 + \rho \lVert C \rVert^2$.

\paragraph{Interaction-lifted variant.}
% code: .../modeling/interaction_lift.py:39 (augment_with_interaction), :62 (step)
% equation_operation: replace the scalar control by the two-column control [v, v*y_t] so the
%   fitted action effect is state-dependent, i.e. the same reminder buys a different amount
%   depending on where the trajectory already is.
The control column is replaced by $\bigl[\,v_t,\ v_t\,y_t\,\bigr]$, so
\eqref{eq:beh-dynamics} becomes
\begin{equation}
  \eta_{t+1} = A\,\eta_t + B_1 v_t + B_2\,\bigl(v_t\,y_t\bigr) + b.
  \label{eq:beh-interaction}
\end{equation}

\paragraph{Note (two ridge implementations).}
This line solves its normal equations directly, while \eqref{eq:stage-two-solve} uses a
pseudoinverse. They are separate implementations in separate packages, deliberately not
shared. Nothing in the paper may describe them as one estimator.

% -------------------------------------------------------------------------------------
% E6. The closed-loop objective, and the budget as tree pruning
% code: persona_drift_control/src/persona_drift/control.py:253 (_current_state),
%       :274 (_remaining_budget), :287 (_planning_steps), :301 (_simulate),
%       :312 (next_u_remind)
% equation_operation: enumerate the full binary action tree of depth equal to the planning
%   horizon, score every branch by the summed predicted readout minus a per-reminder repeat
%   penalty, and emit the first action of the best branch; the budget is enforced by PRUNING
%   -- once the allowance is spent, every deeper node's admissible set collapses to {0}.
% -------------------------------------------------------------------------------------
\begin{align}
  V\bigl(z, a, k, \beta\bigr)
    &= C\,\psi\bigl(f(z,a)\bigr) - \kappa\,\mathbb{1}[a = 1]
     + \max_{a' \in \mathcal{A}(\beta - a)} V\bigl(f(z,a),\, a',\, k-1,\, \beta - a\bigr),
  \label{eq:mpc-value}\\[4pt]
  \mathcal{A}(\beta) &=
  \begin{cases}
    \{0, 1\}, & \beta \ge 1,\\
    \{0\}, & \beta < 1,
  \end{cases}
  \qquad
  u_t = \arg\max_{a \in \{0,1\}} V\bigl(z_t, a, k_t - 1, \beta_t\bigr),
  \label{eq:mpc-budget}
\end{align}
%
\noindent with $V(\cdot,\cdot,0,\cdot)$ terminating at the leading term, $f$ the one-step map
of \eqref{eq:beh-dynamics}, $\kappa$ the repeat penalty, $\beta_t$ the reminders still
unspent, and $k_t = \min(\text{horizon},\ T - t + 1)$.

\paragraph{Note (do not write this as a constrained program).}
\eqref{eq:mpc-budget} is the whole of the constraint mechanism. There is no Lagrangian, no
projection and no relaxation: writing the objective as
$\max \sum_t \hat y_t \ \text{s.t.}\ \sum_t u_t \le \beta$ would misdescribe the executed
controller. At $\beta = 1$ the recursion reduces to a two-way comparison, spending now
against saving the allowance, because choosing $a=1$ forces the entire remaining subtree to
the all-zeros branch.

% -------------------------------------------------------------------------------------
% E7. Marginal benefit of acting, and the threshold it implies
% code: NOT executed anywhere. Derived in
%       docs/experiments/koopman_phaseI_policy_closed_form.md and cross-checked there
%       against control.py:301 (_simulate) on a saved fitted-model JSON.
% equation_operation: evaluate, under the interaction-lifted surrogate, the one-step gain in
%   predicted readout from reminding rather than not, as an affine function of the current
%   readout, and solve its zero crossing for the readout level at which the controller is
%   indifferent.
% -------------------------------------------------------------------------------------
\begin{align}
  m_1(y) &= C\,(I + A)\,\bigl(B_1 + B_2\,y\bigr),
  \qquad
  y^{\star}_{\text{1-step}} = -\,\frac{C(I+A)B_1}{C(I+A)B_2}.
  \label{eq:margin-closed-form}
\end{align}

\paragraph{Note (this is a special case, not the decision rule).}
\eqref{eq:margin-closed-form} equals the controller's actual margin
$V(z,1,\cdot,\cdot) - V(z,0,\cdot,\cdot)$ from \eqref{eq:mpc-value} only when both branches'
optimal second-step action is $0$, which was verified at a single observed grid point. The
numbers differ: $y^{\star}_{\text{1-step}} = 0.8036$ against $y^{\star}_{\text{actual}} =
0.7881$ from the two-step recursion. They are operationally identical here only because the
readout lives on a $\{0, 0.25, 0.5, 0.75, 1.0\}$ grid with no point between them, so both
imply "remind iff $y_{t-1} \le 0.75$". A related derivation is outright wrong and must not
appear: iterating the fixed point $y \leftarrow C(I-A)^{-1}(b + B_1 + B_2 y)$ for the
$v \equiv 1$ steady state diverges, because the augmented control $[1, y]$ is state dependent
and not a constant input. The correct steady state folds the interaction into the system
matrix, $\tilde A = A + B_2 e_1^\top$, before inverting $(I - \tilde A)$.

% -------------------------------------------------------------------------------------
% E8. Finite-horizon controllability matrix and Gramian
% code: src/koopman_ae/core.py:45 (controllability_diagnostics); an independently duplicated,
%       field-identical copy at .../modeling/koopman.py:50, deliberately not shared
% equation_operation: build the finite-horizon controllability matrix by repeated
%   multiplication by A, take its rank and singular values, and separately accumulate the
%   finite-horizon Gramian and its eigenvalues and condition number, alongside the spectrum
%   and spectral radius of A alone.
% -------------------------------------------------------------------------------------
\begin{align}
  \mathcal{C}_H &= \bigl[\,B,\ AB,\ A^2B,\ \dots,\ A^{H-1}B\,\bigr],
  \qquad
  W_H = \sum_{k=0}^{H-1} \bigl(A^k B\bigr)\bigl(A^k B\bigr)^{\!\top},
  \label{eq:ctrb}\\[4pt]
  \rho(A) &= \max_i \lvert \lambda_i(A) \rvert,
  \qquad
  t_{1/2} = \frac{\ln 2}{-\ln \rho(A)}.
  \label{eq:spectral}
\end{align}

\paragraph{Note (reporting discipline).}
Every quantity in \eqref{eq:ctrb} and \eqref{eq:spectral} is a fit diagnostic. Under the
project's reporting rule none of them may carry a cross-arm comparison; they appear only
beside a task-level result. Two further limits were measured rather than assumed. First,
$t_{1/2}$ does not co-sign with whether memory pays: one task has half-life $11.5$ turns and
no measurable gain from delay embedding, another has half-life $0.82$ and a large one. The
quantity that does track it is the older lag's loading in $A$. Second, the impulse-response
norm $\lVert C A^{k-1} B \rVert$ measures the effect still present at step $k$, not the effect
arriving at step $k$; on a synthetic one-step action into an integrator it still reports a
response length of $4$. On any task with a persistent state it can return only one answer, so
it cannot decide whether an action cashes out in one step, and the paper does not use the
action channel's time structure as mechanism evidence.
